\chapter{Introduction}
\label{chap:intro}

\section{Background and Motivation}
\label{sec:background}

The last decade has witnessed an unprecedented convergence of artificial intelligence with nearly every domain of human endeavour. Engineering education has not been immune to this transformation, yet the dominant products in the EdTech space remain either broad-purpose productivity tools (ChatGPT, Gemini, Copilot) not designed for engineering education, or narrow vertical tools (Wolfram Alpha, MATLAB, Overleaf) that address only one aspect of the engineering workflow. The result is a fragmented ecosystem of disconnected point solutions. A typical session today involves a browser for paper search, a notebook for code, a chatbot for questions, and a separate tool for analysis---five different contexts, no shared memory, and a significant cognitive cost at each switch. Sweller's Cognitive Load Theory~\citep{sweller1988cognitive} quantifies this cost: every context switch generates extraneous load that reduces learning efficiency without any corresponding benefit.

Engunity~X was conceived to close this gap. It is a SaaS platform built around four interlocking capabilities---an adaptive retrieval-augmented generation (RAG) pipeline, an autonomous research agent, a secure sandboxed code laboratory, and an adversarial decision-auditing module---all sharing a common memory architecture, authentication layer, and streaming API.

\subsection{The Engineering Education Context}
\label{subsec:edu-context}

Engineering undergraduate programmes in India typically span eight semesters, culminating in a major project. Each project activity maps directly onto one of Engunity's four capabilities:

\begin{enumerate}
    \item \textbf{Literature survey} $\rightarrow$ DeepResearchAgent with OmniRAG retrieval
    \item \textbf{System implementation} $\rightarrow$ Code Lab with AI inline completion and debugging
    \item \textbf{Empirical evaluation} $\rightarrow$ Analytics Workbench with AI-generated insights
    \item \textbf{Critical analysis} $\rightarrow$ Decision Vault with Evidence Quality Score auditing
\end{enumerate}

The alignment is intentional. Engunity was designed by engineering students, for engineering students, with direct knowledge of where the friction points in a final-year project are.

\subsection{The Role of Large Language Models in Technical Assistance}
\label{subsec:llm-role}

Large language models (LLMs) such as GPT-4~\citep{openai2023gpt4} demonstrate remarkable performance on code generation, mathematical reasoning, and scientific question answering. However, raw LLM capability does not translate directly into a reliable engineering assistant: hallucination~\citep{ji2023survey}, stale training knowledge, context blindness, and sycophancy~\citep{perez2022sycophancy} each limit utility. Engunity addresses these failure modes directly: RAG grounds responses in retrieved documents; CRAG web fallback overcomes knowledge staleness; dual-layer memory provides project-specific context; and the Decision Vault's adversarial auditing directly counters sycophancy.

\section{Problem Statement and Objectives}
\label{sec:problem}

\subsection{Core Problem: Tool Fragmentation in Engineering Workflows}
\label{subsec:fragmentation}

The core problem Engunity addresses is \textbf{tool fragmentation}: the absence of a unified environment in which the cognitive activities of engineering work---retrieve, reason, code, analyse, decide---can be performed with shared context and persistent state. Every manual copy-paste between a chatbot and an IDE, every re-explanation of context to a new tool, every search for a paper already reviewed---these are extraneous load events with no learning value.

\subsection{Technical Challenges}
\label{subsec:challenges}

The project addresses five concrete sub-problems: (1) shallow AI reasoning on complex multi-hop queries; (2) code generation without integrated execution and error correction; (3) stateless interaction that loses user context between sessions; (4) multimodal gaps preventing engineering image reasoning; and (5) sycophantic AI that validates incorrect reasoning rather than challenging it.

\subsection{Project Objectives}
\label{subsec:objectives}

In response to these challenges, this project established seven measurable objectives:

\begin{enumerate}
    \item \textbf{OmniRAG Pipeline}: A four-tier complexity-adaptive retrieval system, achieving at least 85\% accuracy on multi-hop queries.
    \item \textbf{DeepResearchAgent}: An autonomous agent with iterative, multi-phase research and configurable depth tiers.
    \item \textbf{Multimodal Input Layer}: A perception stack combining CLIP, YOLOv8n, and EasyOCR, integrated into the RAG pipeline.
    \item \textbf{Secure Code Laboratory}: A containerized execution environment with Monaco Editor IDE, WebSocket terminal, DAP debugger, and AI-assisted error correction.
    \item \textbf{Dual-Layer Memory}: Short-term session summaries combined with long-term per-user profiles persisting across sessions.
    \item \textbf{Decision Vault}: An adversarial reasoning module with a calibrated Evidence Quality Score (EQS) and counter-argument generation.
    \item \textbf{Performance}: Sub-500\,ms P95 end-to-end latency under 500 concurrent users.
\end{enumerate}

\section{Scope of the Project}
\label{sec:scope}

Table~\ref{tab:scope} defines the boundary conditions of this project.

\begin{table}[ht]
    \centering
    \caption{Project scope boundaries.}
    \label{tab:scope}
    \renewcommand{\arraystretch}{1.3}
    \begin{tabular}{p{5.5cm} p{5.5cm}}
        \toprule
        \textbf{In Scope} & \textbf{Out of Scope} \\
        \midrule
        Web-based SaaS platform (browser client) & Native mobile applications \\
        Python, JavaScript, Go, Rust execution & Java, C\#, R execution \\
        Supabase + PostgreSQL authentication & LDAP / enterprise SSO \\
        Single-instance Docker deployment & Kubernetes auto-scaling \\
        English-language queries & Multi-lingual support \\
        CS/Engineering domain documents & Medical, legal domain specialisation \\
        FAISS in-process vector store & Distributed vector database (Qdrant) \\
        Groq API cloud inference & Fine-tuned on-premise LLM hosting \\
        PDF, DOCX, TXT, CSV ingestion & Audio and video ingestion \\
        \bottomrule
    \end{tabular}
\end{table}

\section{Comparison with Existing Tools}
\label{sec:comparison}

Table~\ref{tab:tool-comparison} situates Engunity~X relative to the most widely used AI-assisted tools in engineering education. No existing tool covers all six capability dimensions; Engunity~X is the only platform that combines all of them in a single integrated environment with shared context and memory.

\begin{table}[ht]
    \centering
    \caption{Feature comparison of Engunity~X against existing AI-assisted engineering tools.}
    \label{tab:tool-comparison}
    \renewcommand{\arraystretch}{1.3}
    \resizebox{\linewidth}{!}{%
    \begin{tabular}{lcccccc}
        \toprule
        \textbf{Tool} & \textbf{Doc RAG} & \textbf{Code Exec} & \textbf{Agent} & \textbf{Multimodal} & \textbf{Memory} & \textbf{Adversarial} \\
        \midrule
        ChatGPT 4o     & \checkmark & Limited & \checkmark & \checkmark & Session   & $\times$ \\
        GitHub Copilot & $\times$   & \checkmark & $\times$  & $\times$   & File      & $\times$ \\
        Gemini 1.5     & \checkmark & Limited & \checkmark & \checkmark & Session   & $\times$ \\
        Perplexity AI  & \checkmark & $\times$  & Limited   & $\times$   & $\times$  & $\times$ \\
        Wolfram Alpha  & $\times$   & \checkmark & $\times$ & $\times$   & $\times$  & $\times$ \\
        Replit AI      & $\times$   & \checkmark & $\times$ & $\times$   & Project   & $\times$ \\
        \textbf{Engunity~X} & \checkmark & \checkmark & \checkmark & \checkmark & Long-term & \checkmark \\
        \bottomrule
    \end{tabular}}
\end{table}

\section{Principal Contributions}
\label{sec:contributions}

The principal novel contributions of this project are:

\begin{itemize}
    \item \textbf{OmniRAG}: A four-tier complexity-adaptive retrieval pipeline with a DistilBERT classifier trained on 18,500 labelled queries, achieving 88\% accuracy on multi-hop queries versus 64\% for a standard RAG baseline---a 24-percentage-point improvement.
    \item \textbf{DeepResearchAgent}: A gap-analysis-driven iterative research loop with four configurable depth tiers, integrating vector, graph, and live web retrieval in a single coherent agent.
    \item \textbf{Secure Execution with AI Feedback Loop}: A Docker-sandboxed code execution environment where execution errors are automatically fed back to the LLM for autonomous correction.
    \item \textbf{Evidence Quality Score (EQS)}: A calibrated adversarial reasoning metric on a four-point scale, correctly flagging 78\% of reasoning errors in evaluation.
    \item \textbf{Integrated Dual-Layer Memory}: A persistence system combining LLM-compressed session summaries with a long-term per-user profile spanning sessions.
\end{itemize}

\section{Organisation of the Report}
\label{sec:organisation}

Chapter~\ref{chap:litrev} reviews RAG architectures, autonomous agents, multimodal perception, secure sandboxing, and pedagogical positioning, closing with a gap analysis. Chapter~\ref{chap:design} provides a full technical specification of all five subsystems. Chapter~\ref{chap:results} presents the empirical evaluation. Chapter~\ref{chap:conclusions} summarises contributions, evaluates objectives, discusses limitations, and proposes future work. The Appendices provide API contracts, database schema DDL, environment variable reference, a glossary, and a list of abbreviations.
